\textbf{\textit{Soal. }}Let $ABC$ be an acute triangle with circumcenter $O$, and let $D$, $E$, and $F$ be the feet of altitudes from $A$, $B$, and $C$ to sides $BC$, $CA$, and $AB$, respectively. Denote by $P$ the intersection of the tangents to the circumcircle of $ABC$ at $B$ and $C$. The line through $P$ perpendicular to $EF$ meets $AD$ at $Q$, and let $R$ be the foot of the perpendicular from $A$ to $EF$. Prove that $DR$ and $OQ$ are parallel.

\begin{proof}[\textbf{Solusi.} (Sunday, 9 October 2022) ] 
Misalkan $M$ adalah titik tengah $BC$. Misalkan pula $H$ adalah titik tinggi $\triangle ABC$. 

\begin{lemma}
    $\triangle ARE \sim \triangle ADB$ dan $O,R,A$ segaris
    \begin{proof}
        Karena $\dangle AER = \dangle AEF = \dangle CEF = \dangle CBF = \dangle DBA$ dan $\dangle ARE = \dangle BDA = 90^\circ$, maka $\triangle ARE \sim \triangle ADB$ sehingga $\dangle RAE = \dangle BAD$. Namun, karena \textit{well-known} bahwa $\dangle CAO = 90^\circ - \dangle B = \dangle DAB$, maka kita punya $\dangle RAE = \dangle OAC$ yang menyebabkan $O,R,A$ segaris 
    \end{proof}
\end{lemma}
\begin{lemma}
    $OAQP$ jajargenjang
    \begin{proof}
        Jelas bahwa $O,M,P$ segaris dan $OP \perp BC$. Tetapi, karena $AQ \perp BC$ juga, maka $OP \parallel AQ$. Lalu, karena $O,R,A$ segaris, maka $OA \perp FE$. Namun, karena $QP \perp FE$ juga, maka $OA \parallel QP$. Dari sini didapatkan bahwa $OAQP$ jajargenjang.
    \end{proof}
\end{lemma}

Dikarenakan $OAQP$ jajargenjang, maka $OP=OB$. Lalu, $OB=OA$ yang merupakan jari-jari lingkaran luar $(ABC)$. Selanjutnya, karena $\triangle ARE \sim \triangle ADB$ dan karena jelas bahwa $\triangle OBM \sim \triangle OPB$, maka 
\begin{align*}
    \dfrac{AR}{AD} = \dfrac{AE}{AB} = \cos \angle A = \dfrac{OM}{OB} = \dfrac{OB}{OP} = \dfrac{OA}{AQ}
\end{align*}
karena $\dangle RAD = \dangle OAQ$ juga, maka ditemukan bahwa $\triangle RAD \sim \triangle OAQ$. Ini berarti $\dangle RDA = \dangle OQA \implies DR \parallel OQ$. Terbukti.

\begin{center}
% Setting up the TikZ picture
\begin{tikzpicture}[scale=2.5, dot/.style={circle, fill, inner sep=0.5pt, outer sep=0pt}]

    % Defining the unit circle radius and points A, B, C using polar coordinates
    \tkzDefPoint(125:1){A}
    \tkzDefPoint(-150:1){B}
    \tkzDefPoint(-30:1){C}
    
    % Defining origin O
    \tkzDefPoint(0,0){O}
    
    % Computing point H as the sum of vectors A, B, C
    \tkzDefTriangleCenter[ortho](A,B,C) \tkzGetPoint{H}
    
    % Computing midpoint M of B and C
    \tkzDefMidPoint(B,C) \tkzGetPoint{M}
    
    % Computing feet of perpendiculars (projections)
    % D is the foot of perpendicular from A to line BC
    \tkzDefPointBy[projection=onto B--C](A) \tkzGetPoint{D}
    % E is the foot of perpendicular from B to line CA
    \tkzDefPointBy[projection=onto C--A](B) \tkzGetPoint{E}
    % F is the foot of perpendicular from C to line AB
    \tkzDefPointBy[projection=onto A--B](C) \tkzGetPoint{F}
    
    % P is the intersection of tangent from B and C to the circle
    \tkzDefLine[tangent at=B](O) \tkzGetPoint{refB}
    \tkzDefLine[tangent at=C](O) \tkzGetPoint{refC}
    \tkzInterLL(B,refB)(C,refC) \tkzGetPoint{P}
    \tkzDrawSegment(B,P)
    \tkzDrawSegment(C,P)

    % Computing point Q = A + P - O
    \tkzDefPointBy[translation=from O to P](A) \tkzGetPoint{Q}
    
    % Computing point R as the intersection of lines A-O and E-F
    \tkzInterLL(A,O)(E,F) \tkzGetPoint{R}
    
    % Computing point S as the intersection of lines E-F and P-Q
    \tkzInterLL(E,F)(P,Q) \tkzGetPoint{S}
    
    % Drawing the unit circle (dotted)
    \tkzDrawCircle[dashed](O,A)
    
    % Drawing triangle A-B-C
    \tkzDrawPolygon(A,B,C)
    
    % Drawing lines P-B and P-C
    \tkzDrawSegment(P,B)
    \tkzDrawSegment(P,C)
    
    % Drawing dashed grey lines F-S and Q-S
    \tkzDrawSegment[dashed, gray](F,S)
    \tkzDrawSegment[dashed, gray](Q,S)
    
    % Drawing blue path A-Q-P-O-A
    \tkzDrawPolygon[blue](A,Q,P,O)
    
    % Drawing lines B-E, C-F, E-F
    \tkzDrawSegment(B,E)
    \tkzDrawSegment(C,F)
    \tkzDrawSegment(E,F)
    
    % Placing labeled dots
    \tkzDrawPoint[dot, fill=black, size=3](A) \tkzLabelPoint[above](A){$A$}
    \tkzDrawPoint[dot, fill=black, size=3](B) \tkzLabelPoint[below left](B){$B$}
    \tkzDrawPoint[dot, fill=black, size=3](C) \tkzLabelPoint[below right](C){$C$}
    \tkzDrawPoint[dot, fill=black, size=3](P) \tkzLabelPoint[below right](P){$P$}
    \tkzDrawPoint[dot, fill=black, size=3](O) \tkzLabelPoint[below](O){$O$}
    \tkzDrawPoint[dot, fill=black, size=3](H) \tkzLabelPoint[above](H){$H$}
    \tkzDrawPoint[dot, fill=black, size=3](D) \tkzLabelPoint[below right](D){$D$}
    \tkzDrawPoint[dot, fill=black, size=3](E) \tkzLabelPoint[above right](E){$E$}
    \tkzDrawPoint[dot, fill=black, size=3](F) \tkzLabelPoint[left](F){$F$}
    \tkzDrawPoint[dot, fill=black, size=3](R) \tkzLabelPoint[left](R){$R$}
    \tkzDrawPoint[dot, fill=black, size=3](Q) \tkzLabelPoint[left](Q){$Q$}
    \tkzDrawPoint[dot, fill=black, size=3](S)
    \tkzDrawPoint[dot, fill=black, size=3](M) \tkzLabelPoint[below right](M){$M$}
    
\end{tikzpicture}
\end{center}
\end{proof}